\documentclass{amsart}
\usepackage[UTF8,fontset=windows]{ctex}
\usepackage[a4paper,margin=24mm]{geometry}
% For dealing with references we use the comment environment
\usepackage{verbatim}
\newenvironment{reference}{\comment}{\endcomment}
%\newenvironment{reference}{}{}
\newenvironment{slogan}{\comment}{\endcomment}
\newenvironment{history}{\comment}{\endcomment}

% For commutative diagrams we use Xy-pic
\usepackage[all]{xy}

% We use 2cell for 2-commutative diagrams.
\xyoption{2cell}
\UseAllTwocells

% We use multicol for the list of chapters between chapters
\usepackage{multicol}

% This is generally recommended for better output
\usepackage{lmodern}
\usepackage[T1]{fontenc}

% For cross-file-references

% Package for hypertext links:
\usepackage{hyperref}

% For any local file, say "hello.tex" you want to link to please

% Theorem environments.
%
\theoremstyle{plain}
\newtheorem{theorem}[subsection]{Theorem}
\newtheorem{proposition}[subsection]{Proposition}
\newtheorem{lemma}[subsection]{Lemma}

\theoremstyle{definition}
\newtheorem{definition}[subsection]{Definition}
\newtheorem{example}[subsection]{Example}
\newtheorem{exercise}[subsection]{Exercise}
\newtheorem{situation}[subsection]{Situation}

\theoremstyle{remark}
\newtheorem{remark}[subsection]{Remark}
\newtheorem{remarks}[subsection]{Remarks}

\numberwithin{equation}{subsection}

% Macros
%
\def\lim{\mathop{\mathrm{lim}}\nolimits}
\def\colim{\mathop{\mathrm{colim}}\nolimits}
\def\Spec{\mathop{\mathrm{Spec}}}
\def\Hom{\mathop{\mathrm{Hom}}\nolimits}
\def\Ext{\mathop{\mathrm{Ext}}\nolimits}
\def\SheafHom{\mathop{\mathcal{H}\!\mathit{om}}\nolimits}
\def\SheafExt{\mathop{\mathcal{E}\!\mathit{xt}}\nolimits}
\def\Sch{\mathit{Sch}}
\def\Mor{\mathop{\mathrm{Mor}}\nolimits}
\def\Ob{\mathop{\mathrm{Ob}}\nolimits}
\def\Sh{\mathop{\mathit{Sh}}\nolimits}
\def\NL{\mathop{N\!L}\nolimits}
\def\CH{\mathop{\mathrm{CH}}\nolimits}
\def\proetale{{pro\text{-}\acute{e}tale}}
\def\etale{{\acute{e}tale}}
\def\QCoh{\mathit{QCoh}}
\def\Ker{\mathop{\mathrm{Ker}}}
\def\Im{\mathop{\mathrm{Im}}}
\def\Coker{\mathop{\mathrm{Coker}}}
\def\Coim{\mathop{\mathrm{Coim}}}

% Boxtimes
%
\DeclareMathSymbol{\boxtimes}{\mathbin}{AMSa}{"02}

%
% Macros for moduli stacks/spaces
%
\def\QCohstack{\mathcal{QC}\!\mathit{oh}}
\def\Cohstack{\mathcal{C}\!\mathit{oh}}
\def\Spacesstack{\mathcal{S}\!\mathit{paces}}
\def\Quotfunctor{\mathrm{Quot}}
\def\Hilbfunctor{\mathrm{Hilb}}
\def\Curvesstack{\mathcal{C}\!\mathit{urves}}
\def\Polarizedstack{\mathcal{P}\!\mathit{olarized}}
\def\Complexesstack{\mathcal{C}\!\mathit{omplexes}}
% \Pic is the operator that assigns to X its picard group, usage \Pic(X)
% \Picardstack_{X/B} denotes the Picard stack of X over B
% \Picardfunctor_{X/B} denotes the Picard functor of X over B
\def\Pic{\mathop{\mathrm{Pic}}\nolimits}
\def\Picardstack{\mathcal{P}\!\mathit{ic}}
\def\Picardfunctor{\mathrm{Pic}}
\def\Deformationcategory{\mathcal{D}\!\mathit{ef}}

\hypersetup{hidelinks}
\widowpenalty=10000\clubpenalty=10000\displaywidowpenalty=10000
\begin{document}
\title[Stacks Project, Tag 01PJ]{QCQS 概形上的拟凝聚层有限呈示逼近}
\maketitle
\noindent\textbf{Counting scope.} The primary problem is Tag 01PJ, marked below. All preceding numbered statements are supporting author context and are not counted as additional problems. This is a local pilot with a temporary ID, not production admission.


\noindent Copyright (C) 2005--2025 Johan de Jong. Stacks Project authors. Source: \href{https://stacks.math.columbia.edu/tag/01PJ}{Tag 01PJ}.

\noindent Difficulty H2: 本题在拟紧拟分离概形上把任意拟凝聚层实现为有限呈示层的滤过余极限。证明构造所有有限呈示层到目标的范畴，用直和与余核验证滤过性，并借局部核及截面的全局延拓证明余极限映射在茎上的单射和满射。有限呈示延拓是输入，索引范畴及茎级比较仍由本证明承担，判H2。.

\noindent History: excerpt from revision local cached snapshot; exact chapter hashes; revision unknown, properties.tex, lines 3039--3190. Only reference presentation was adapted. Licensed under GNU FDL 1.2 or later; no invariant sections or cover texts. See accompanying COPYING.

\begin{lemma}
\label{lemma-extend}
Let $X$ be a quasi-compact and quasi-separated scheme.
Let $U \subset X$ be a quasi-compact open.
Let $\mathcal{F}$ be a quasi-coherent $\mathcal{O}_X$-module.
Let $\mathcal{G} \subset \mathcal{F}|_U$ be a quasi-coherent
$\mathcal{O}_U$-submodule which is of finite type. Then
there exists a quasi-coherent submodule $\mathcal{G}' \subset \mathcal{F}$
which is of finite type such that $\mathcal{G}'|_U = \mathcal{G}$.
\end{lemma}

\begin{proof}
Let $n$ be the minimal number of affine opens $U_i \subset X$,
$i = 1, \ldots , n$ such that $X = U \cup \bigcup U_i$.
(Here we use that $X$ is quasi-compact.) Suppose
we can prove the lemma for the case $n = 1$. Then we can successively
extend $\mathcal{G}$
to a $\mathcal{G}_1$ over $U \cup U_1$
to a $\mathcal{G}_2$ over $U \cup U_1 \cup U_2$
to a $\mathcal{G}_3$ over $U \cup U_1 \cup U_2 \cup U_3$,
and so on.
Thus we reduce to the case $n = 1$.

\medskip\noindent
Thus we may assume that $X = U \cup V$ with $V$ affine.
Since $X$ is quasi-separated and $U$, $V$ are quasi-compact open,
we see that $U \cap V$ is a quasi-compact open. It suffices to prove the
lemma for the system $(V, U \cap V, \mathcal{F}|_V, \mathcal{G}|_{U \cap V})$
since we can glue the resulting sheaf $\mathcal{G}'$ over $V$
to the given sheaf $\mathcal{G}$ over $U$ along the common value
over $U \cap V$.
Thus we reduce to the case where $X$ is affine.

\medskip\noindent
Assume $X = \Spec(R)$. Write $\mathcal{F} = \widetilde M$
for some $R$-module $M$. By Lemma \href{https://stacks.math.columbia.edu/tag/01PE}{lemma extend trivial (Tag 01PE)} above we may
find a quasi-coherent subsheaf $\mathcal{H} \subset \mathcal{F}$
which restricts to $\mathcal{G}$ over $U$.
Write $\mathcal{H} = \widetilde N$ for some $R$-module $N$.
For every $u \in U$ there exists an $f \in R$ such that
$u \in D(f) \subset U$ and such that $N_f$ is finitely generated,
see Lemma \href{https://stacks.math.columbia.edu/tag/01PB}{lemma finite type module (Tag 01PB)}.
Since $U$ is quasi-compact we can cover it by finitely
many $D(f_i)$ such that $N_{f_i}$ is generated by
finitely many elements, say $x_{i, 1}/f_i^N, \ldots, x_{i, r_i}/f_i^N$.
Let $N' \subset N$ be the submodule generated by the elements
$x_{i, j}$. Then the subsheaf
$\mathcal{G}' = \widetilde{N'} \subset \mathcal{H} \subset \mathcal{F}$
works.
\end{proof}


\begin{lemma}
\label{lemma-extend-finite-presentation}
Let $X$ be a quasi-compact and quasi-separated scheme.
Let $\mathcal{F}$ be a quasi-coherent $\mathcal{O}_X$-module.
Let $U \subset X$ be a quasi-compact open.
Let $\mathcal{G}$ be an $\mathcal{O}_U$-module which is of finite presentation.
Let $\varphi : \mathcal{G} \to \mathcal{F}|_U$ be a morphism of
$\mathcal{O}_U$-modules.
Then there exists an $\mathcal{O}_X$-module
$\mathcal{G}'$ of finite presentation, and a morphism
of $\mathcal{O}_X$-modules $\varphi' : \mathcal{G}' \to \mathcal{F}$
such that $\mathcal{G}'|_U = \mathcal{G}$ and such that
$\varphi'|_U = \varphi$.
\end{lemma}

\begin{proof}
The beginning of the proof is a repeat of the beginning of the
proof of Lemma \ref{lemma-extend}. We write it out carefully anyway.

\medskip\noindent
Let $n$ be the minimal number of affine opens $U_i \subset X$,
$i = 1, \ldots , n$ such that $X = U \cup \bigcup U_i$.
(Here we use that $X$ is quasi-compact.) Suppose
we can prove the lemma for the case $n = 1$. Then we can successively
extend the pair $(\mathcal{G}, \varphi)$
to a pair $(\mathcal{G}_1, \varphi_1)$ over $U \cup U_1$
to a pair $(\mathcal{G}_2, \varphi_2)$ over $U \cup U_1 \cup U_2$
to a pair $(\mathcal{G}_3, \varphi_3)$ over $U \cup U_1 \cup U_2 \cup U_3$,
and so on.
Thus we reduce to the case $n = 1$.

\medskip\noindent
Thus we may assume that $X = U \cup V$ with $V$ affine.
Since $X$ is quasi-separated and $U$ quasi-compact,
we see that $U \cap V \subset V$ is quasi-compact.
Suppose we prove the lemma for the system
$(V, U \cap V, \mathcal{F}|_V, \mathcal{G}|_{U \cap V}, \varphi|_{U \cap V})$
thereby producing $(\mathcal{G}', \varphi')$ over $V$.
Then we can glue $\mathcal{G}'$ over $V$ to the given sheaf $\mathcal{G}$
over $U$ along the common value over $U \cap V$, and similarly we can glue
the map $\varphi'$ to the map $\varphi$ along the common value over
$U \cap V$. Thus we reduce to the case where $X$ is affine.

\medskip\noindent
Assume $X = \Spec(R)$.
By Lemma \href{https://stacks.math.columbia.edu/tag/01PE}{lemma extend trivial (Tag 01PE)} above we may
find a quasi-coherent sheaf $\mathcal{H}$ with
a map $\psi : \mathcal{H} \to \mathcal{F}$ over $X$
which restricts to $\mathcal{G}$ and $\varphi$ over $U$.
By Lemma \ref{lemma-extend} we can find a finite type
quasi-coherent $\mathcal{O}_X$-submodule
$\mathcal{H}' \subset \mathcal{H}$
such that $\mathcal{H}'|_U = \mathcal{G}$. Thus after
replacing $\mathcal{H}$ by $\mathcal{H}'$
and $\psi$ by the restriction of $\psi$ to $\mathcal{H}'$
we may assume that $\mathcal{H}$ is of finite type.
By Lemma \href{https://stacks.math.columbia.edu/tag/01PC}{lemma finite presentation module (Tag 01PC)}
we conclude that $\mathcal{H} = \widetilde{N}$ with
$N$ a finitely generated $R$-module. Hence there exists a surjection
as in the following short exact sequence of
quasi-coherent $\mathcal{O}_X$-modules
$$
0 \to \mathcal{K} \to \mathcal{O}_X^{\oplus n} \to \mathcal{H} \to 0
$$
where $\mathcal{K}$ is defined as the kernel.
Since $\mathcal{G}$ is of finite presentation and
$\mathcal{H}|_U = \mathcal{G}$ by
Modules, Lemma
\href{https://stacks.math.columbia.edu/tag/01BP}{lemma kernel surjection finite free onto finite presentation (Tag 01BP)}
the restriction $\mathcal{K}|_U$ is
an $\mathcal{O}_U$-module of finite type. Hence by Lemma \ref{lemma-extend}
again we see that there exists a finite type quasi-coherent
$\mathcal{O}_X$-submodule $\mathcal{K}' \subset \mathcal{K}$ such
that $\mathcal{K}'|_U = \mathcal{K}|_U$. The solution to the problem
posed in the lemma is to set
$$
\mathcal{G}' = \mathcal{O}_X^{\oplus n}/\mathcal{K}'
$$
which is clearly of finite presentation and restricts to give $\mathcal{G}$
on $U$ with $\varphi'$ equal to the composition
$$
\mathcal{G}' = \mathcal{O}_X^{\oplus n}/\mathcal{K}'
\to \mathcal{O}_X^{\oplus n}/\mathcal{K} = \mathcal{H} \xrightarrow{\psi}
\mathcal{F}.
$$
This finishes the proof of the lemma.
\end{proof}


\begin{definition}
\label{categories-definition-directed}
We say that a diagram $M : \mathcal{I} \to \mathcal{C}$ is {\it directed},
or {\it filtered} if the following conditions hold:
\begin{enumerate}
\item the category $\mathcal{I}$ has at least one object,
\item for every pair of objects $x, y$ of $\mathcal{I}$
there exist an object $z$ and morphisms $x \to z$,
$y \to z$, and
\item for every pair of objects $x, y$ of $\mathcal{I}$
and every pair of morphisms $a, b : x \to y$ of $\mathcal{I}$
there exists a morphism $c : y \to z$ of $\mathcal{I}$
such that $M(c \circ a) = M(c \circ b)$ as morphisms in $\mathcal{C}$.
\end{enumerate}
We say that an index category $\mathcal{I}$ is {\it directed}, or
{\it filtered} if $\text{id} : \mathcal{I} \to \mathcal{I}$ is filtered
(in other words you erase the $M$ in part (3) above).
\end{definition}

\section*{Primary problem: Tag 01PJ}
\begin{lemma}
\label{lemma-directed-colimit-diagram-finite-presentation}
\begin{slogan}
Quasi-coherent modules on quasi-compact and quasi-separated schemes
are filtered colimits of finitely presented modules.
\end{slogan}
Let $X$ be a scheme. Assume $X$ is quasi-compact and quasi-separated.
Let $\mathcal{F}$ be a quasi-coherent $\mathcal{O}_X$-module.
There exist
\begin{enumerate}
\item a filtered index category $\mathcal{I}$ (see
Categories, Definition \ref{categories-definition-directed}),
\item a diagram $\mathcal{I} \to \textit{Mod}(\mathcal{O}_X)$ (see
Categories, Section \href{https://stacks.math.columbia.edu/tag/002D}{section limits (Tag 002D)}),
$i \mapsto \mathcal{F}_i$,
\item morphisms of $\mathcal{O}_X$-modules
$\varphi_i : \mathcal{F}_i \to \mathcal{F}$
\end{enumerate}
such that each $\mathcal{F}_i$ is of finite presentation
and such that the morphisms $\varphi_i$ induce an isomorphism
$$
\colim_i \mathcal{F}_i
=
\mathcal{F}.
$$
\end{lemma}

\begin{proof}
Choose a set $I$ and for each $i \in I$ an $\mathcal{O}_X$-module
of finite presentation and a homomorphism of $\mathcal{O}_X$-modules
$\varphi_i : \mathcal{F}_i \to \mathcal{F}$ with the following
property: For any $\psi : \mathcal{G} \to \mathcal{F}$ with $\mathcal{G}$
of finite presentation there is an $i \in I$ such that there exists
an isomorphism $\alpha : \mathcal{F}_i \to \mathcal{G}$ with
$\varphi_i = \psi \circ \alpha$. It is clear from
Modules, Lemma \href{https://stacks.math.columbia.edu/tag/01BC}{lemma set isomorphism classes finite type modules (Tag 01BC)}
that such a set exists (see also its proof).
We denote $\mathcal{I}$ the category
with $\Ob(\mathcal{I}) = I$ and given $i, i' \in I$
we set
$$
\Mor_\mathcal{I}(i, i') =
\{\alpha : \mathcal{F}_i \to \mathcal{F}_{i'} \mid
\alpha \circ \varphi_{i'} = \varphi_i
\}.
$$
We claim that $\mathcal{I}$ is a filtered category and that
$\mathcal{F} = \colim_i \mathcal{F}_i$.

\medskip\noindent
Let $i, i' \in I$. Then we can consider the morphism
$$
\mathcal{F}_i \oplus \mathcal{F}_{i'} \longrightarrow \mathcal{F}
$$
which is the direct sum of $\varphi_i$ and $\varphi_{i'}$.
Since a direct sum of finitely presented $\mathcal{O}_X$-modules
is finitely presented we see that there exists some $i'' \in I$
such that $\varphi_{i''} : \mathcal{F}_{i''} \to \mathcal{F}$
is isomorphic to the displayed arrow towards $\mathcal{F}$ above.
Since there are commutative diagrams
$$
\xymatrix{
\mathcal{F}_i \ar[r] \ar[d] & \mathcal{F} \ar@{=}[d] \\
\mathcal{F}_i \oplus \mathcal{F}_{i'} \ar[r] & \mathcal{F}
}
\quad
\text{and}
\quad
\xymatrix{
\mathcal{F}_{i'} \ar[r] \ar[d] & \mathcal{F} \ar@{=}[d] \\
\mathcal{F}_i \oplus \mathcal{F}_{i'} \ar[r] & \mathcal{F}
}
$$
we see that there are morphisms $i \to i''$ and $i' \to i''$
in $\mathcal{I}$. Next, suppose that we have $i, i' \in I$ and
morphisms $\alpha, \beta : i \to i'$ (corresponding to $\mathcal{O}_X$-module
maps $\alpha, \beta : \mathcal{F}_i \to \mathcal{F}_{i'}$).
In this case consider the coequalizer
$$
\mathcal{G} =
\Coker(
\mathcal{F}_i \xrightarrow{\alpha - \beta} \mathcal{F}_{i'}
)
$$
Note that $\mathcal{G}$ is an $\mathcal{O}_X$-module of finite presentation.
Since by definition of morphisms in the category $\mathcal{I}$
we have $\varphi_{i'} \circ \alpha = \varphi_{i'} \circ \beta$
we see that we get an induced map $\psi : \mathcal{G} \to \mathcal{F}$.
Hence again the pair $(\mathcal{G}, \psi)$ is isomorphic to
the pair $(\mathcal{F}_{i''}, \varphi_{i''})$ for some $i''$.
Hence we see that there exists a morphism $i' \to i''$ in
$\mathcal{I}$ which equalizes $\alpha$ and $\beta$. Thus we have
shown that the category $\mathcal{I}$ is filtered.

\medskip\noindent
We still have to show that the colimit of the diagram is $\mathcal{F}$.
By definition of the colimit, and by our definition of the category
$\mathcal{I}$ there is a canonical map
$$
\varphi :
\colim_i \mathcal{F}_i
\longrightarrow
\mathcal{F}.
$$
Pick $x \in X$. Let us show that $\varphi_x$ is an isomorphism.
Recall that
$$
(\colim_i \mathcal{F}_i)_x
=
\colim_i \mathcal{F}_{i, x},
$$
see
Sheaves, Section \href{https://stacks.math.columbia.edu/tag/009E}{section limits sheaves (Tag 009E)}.
First we show that the map $\varphi_x$ is injective.
Suppose that $s \in \mathcal{F}_{i, x}$ is an element
such that $s$ maps to zero in $\mathcal{F}_x$. Then there exists
a quasi-compact open $U$ such that $s$ comes from $s \in \mathcal{F}_i(U)$
and such that $\varphi_i(s) = 0$ in $\mathcal{F}(U)$.
By Lemma \ref{lemma-extend}
we can find a finite type quasi-coherent subsheaf
$\mathcal{K} \subset \Ker(\varphi_i)$ which restricts to
the quasi-coherent $\mathcal{O}_U$-submodule of $\mathcal{F}_i$
generated by $s$:
$\mathcal{K}|_U = \mathcal{O}_U\cdot s \subset \mathcal{F}_i|_U$.
Clearly, $\mathcal{F}_i/\mathcal{K}$ is of finite presentation and
the map $\varphi_i$ factors through the quotient map
$\mathcal{F}_i \to \mathcal{F}_i/\mathcal{K}$. Hence we can find
an $i' \in I$ and a morphism $\alpha : \mathcal{F}_i \to \mathcal{F}_{i'}$
in $\mathcal{I}$ which can be identified with the quotient map
$\mathcal{F}_i \to \mathcal{F}_i/\mathcal{K}$. Then it follows
that the section $s$ maps to zero in $\mathcal{F}_{i'}(U)$ and
in particular in
$(\colim_i \mathcal{F}_i)_x =
\colim_i \mathcal{F}_{i, x}$.
The injectivity follows.
Finally, we show that the map $\varphi_x$ is surjective.
Pick $s \in \mathcal{F}_x$. Choose a quasi-compact open neighbourhood
$U \subset X$ of $x$ such that $s$ corresponds to a section
$s \in \mathcal{F}(U)$. Consider the map
$s : \mathcal{O}_U \to \mathcal{F}$ (multiplication by $s$).
By Lemma \ref{lemma-extend-finite-presentation}
there exists an $\mathcal{O}_X$-module $\mathcal{G}$
of finite presentation and an $\mathcal{O}_X$-module map
$\mathcal{G} \to \mathcal{F}$ such that $\mathcal{G}|_U \to \mathcal{F}|_U$
is identified with
$s : \mathcal{O}_U \to \mathcal{F}$.
Again by definition of $\mathcal{I}$ there exists an $i \in I$
such that $\mathcal{G} \to \mathcal{F}$ is isomorphic to
$\varphi_i : \mathcal{F}_i \to \mathcal{F}$. Clearly there exists
a section $s' \in \mathcal{F}_i(U)$ mapping to $s \in \mathcal{F}(U)$.
This proves surjectivity and the proof of the lemma is complete.
\end{proof}
\end{document}
